% الجزء: sets-functions-relations
% الفصل: sets
% القسم: subsets

\documentclass[../../../include/open-logic-section]{subfiles}

\begin{document}

\olfileid[ar]{sfr}{set}{sub}
\olsection{المجموعات الجزئية ومجموعات القوى}

\begin{explain}
من وجوه النظر في مجموعتين أن نسأل: أكل ما في إحداهما في الأخرى؟
وهذا وجه كثير الورود في المقارنة بين المجموعات. فلأهميته نخص الحالة
التي يكون فيها \emph{كل ما في مجموعة موجودًا في الأخرى أيضًا} برمز.
\end{explain}

\begin{defn}[المجموعة الجزئية]
المجموعة ${\OLId{latin-looped}{A}}$ \emph{مجموعة جزئية} من~${\OLId{latin-looped}{B}}$ متى كان كل !!{element} من ${\OLId{latin-looped}{A}}$
أيضًا !!a{element} من~${\OLId{latin-looped}{B}}$، ونكتب لذلك ${\OLId{latin-looped}{A}} \subseteq {\OLId{latin-looped}{B}}$.
فإن لم تكن ${\OLId{latin-looped}{A}}$ مجموعة جزئية من~${\OLId{latin-looped}{B}}$ كتبنا ${\OLId{latin-looped}{A}} \not\subseteq {\OLId{latin-looped}{B}}$.
وإن كان ${\OLId{latin-looped}{A}} \subseteq {\OLId{latin-looped}{B}}$ مع ${\OLId{latin-looped}{A}} \neq {\OLId{latin-looped}{B}}$، فكتابته ${\OLId{latin-looped}{A}} \subsetneq {\OLId{latin-looped}{B}}$،
وتسمى ${\OLId{latin-looped}{A}}$ حينئذ \emph{مجموعة جزئية فعلية} من ${\OLId{latin-looped}{B}}$.
\end{defn}

\begin{ex}
كل مجموعة جزئية من نفسها، والخالية $\emptyset$ جزئية من كل مجموعة.
ومن أمثلة ذلك أن الأعداد الطبيعية الزوجية مجموعة جزئية من الأعداد
الطبيعية، وأن $\{ {\OLId{latin-ordinary}{a}}, {\OLId{latin-ordinary}{b}} \} \subseteq \{ {\OLId{latin-ordinary}{a}}, {\OLId{latin-ordinary}{b}}, {\OLId{latin-ordinary}{c}} \}$.
وأما $\{ {\OLId{latin-ordinary}{a}}, {\OLId{latin-ordinary}{b}}, {\OLId{latin-ordinary}{e}}
\}$ فليست مجموعة جزئية من $\{ {\OLId{latin-ordinary}{a}}, {\OLId{latin-ordinary}{b}}, {\OLId{latin-ordinary}{c}} \}$.
\end{ex}

\begin{ex}
ليست العضوية هي الجزئية: فالعدد ${\OLMathNumeral{2}}$ هو !!{element} في مجموعة الأعداد
الصحيحة، وأما مجموعة الأعداد الزوجية فمجموعة جزئية منها. وقد يجتمع
الوصفان لمجموعة بالنسبة إلى أخرى، فيصدق عليها \emph{كلا الأمرين}: كونها
!!a{element} وكونها مجموعة جزئية. ومثاله $\{{\OLMathNumeral{0}}\} \in \{{\OLMathNumeral{0}}, \{{\OLMathNumeral{0}}\}\}$ مع $\{{\OLMathNumeral{0}}\} \subseteq \{{\OLMathNumeral{0}},
\{{\OLMathNumeral{0}}\}\}$.
\end{ex}

معيار الامتدادية لتساوي المجموعات هو أن ${\OLId{latin-looped}{A}} = {\OLId{latin-looped}{B}}$ إذا وفقط إذا كان
كل !!{element} من~${\OLId{latin-looped}{A}}$ أيضًا !!a{element} من~${\OLId{latin-looped}{B}}$، وبالعكس.
فأما شطره الأول فهو بعينه معنى ${\OLId{latin-looped}{A}} \subseteq {\OLId{latin-looped}{B}}$: كل !!{element}
من~${\OLId{latin-looped}{A}}$ هو !!a{element} من~${\OLId{latin-looped}{B}}$. وأما إذا بدّلنا ${\OLId{latin-looped}{A}}$ و${\OLId{latin-looped}{B}}$، كان
${\OLId{latin-looped}{B}} \subseteq {\OLId{latin-looped}{A}}$ معناه أن كل !!{element} من~${\OLId{latin-looped}{B}}$ هو !!a{element}
من~${\OLId{latin-looped}{A}}$؛ وهذا هو شطر «بالعكس». فحاصل الامتدادية أن المجموعتين تتساويان
إذا وفقط إذا كانت كل واحدة منهما جزئية من صاحبتها.

\begin{prop}
شرط ${\OLId{latin-looped}{A}} = {\OLId{latin-looped}{B}}$ اللازم والكافي اجتماع ${\OLId{latin-looped}{A}} \subseteq {\OLId{latin-looped}{B}}$ و${\OLId{latin-looped}{B}} \subseteq {\OLId{latin-looped}{A}}$.
\end{prop}

ولنضع هنا رمزًا آخر ينفع فيما يأتي. فقد قلنا في تعريف كون ${\OLId{latin-looped}{A}}$ جزئية
من~${\OLId{latin-looped}{B}}$: «كل !!{element} من~${\OLId{latin-looped}{A}}$ هو \dots»، وجعلنا موضع «$\dots$»
قولنا «!!a{element} من ${\OLId{latin-looped}{B}}$». وهذا \emph{النمط} من القول كثير الدوران؛
فحسن أن نجعل له اختصارًا صوريًا.

\begin{defn}\ollabel{forallxina}
الرمز $(\forall {\OLId{latin-ordinary}{x}} \in {\OLId{latin-looped}{A}})\phi$ اختصار لـ$\forall {\OLId{latin-ordinary}{x}}({\OLId{latin-ordinary}{x}} \in {\OLId{latin-looped}{A}} \lif
\phi)$، وكذلك $(\exists {\OLId{latin-ordinary}{x}} \in {\OLId{latin-looped}{A}})\phi$ اختصار لـ$\exists {\OLId{latin-ordinary}{x}}({\OLId{latin-ordinary}{x}}
\in {\OLId{latin-looped}{A}} \land \phi)$.
\end{defn}

وعلى هذا الاصطلاح يكون ${\OLId{latin-looped}{A}} \subseteq {\OLId{latin-looped}{B}}$ إذا وفقط إذا كان $(\forall
{\OLId{latin-ordinary}{x}} \in {\OLId{latin-looped}{A}}){\OLId{latin-ordinary}{x}} \in {\OLId{latin-looped}{B}}$.

ولننظر الآن في جمع المجموعات الجزئية كلها من مجموعة معطاة في مجموعة واحدة.

\begin{defn}[مجموعة القوى]
\emph{مجموعة قوى}~${\OLId{latin-looped}{A}}$ هي المجموعة الجامعة لجميع المجموعات الجزئية
من~${\OLId{latin-looped}{A}}$، ورمزها $\Pow{{\OLId{latin-looped}{A}}}$:
  \[
    \Pow{{\OLId{latin-looped}{A}}} = \Setabs{{\OLId{latin-looped}{B}}}{{\OLId{latin-looped}{B}} \subseteq {\OLId{latin-looped}{A}}}
  \]
\end{defn}

\begin{ex}
إذا طلبنا جميع جزئيات $\{ {\OLId{latin-ordinary}{a}}, {\OLId{latin-ordinary}{b}}, {\OLId{latin-ordinary}{c}} \}$ وجدناها:
$\emptyset$، و$\{{\OLId{latin-ordinary}{a}} \}$، و$\{{\OLId{latin-ordinary}{b}}\}$، و$\{{\OLId{latin-ordinary}{c}}\}$، و$\{{\OLId{latin-ordinary}{a}}, {\OLId{latin-ordinary}{b}}\}$، و$\{{\OLId{latin-ordinary}{a}}, {\OLId{latin-ordinary}{c}}\}$، و$\{{\OLId{latin-ordinary}{b}},
{\OLId{latin-ordinary}{c}}\}$، و$\{{\OLId{latin-ordinary}{a}}, {\OLId{latin-ordinary}{b}}, {\OLId{latin-ordinary}{c}}\}$. فالمجموعة التي تجمع هذه كلها هي
$\Pow{\{{\OLId{latin-ordinary}{a}},{\OLId{latin-ordinary}{b}},{\OLId{latin-ordinary}{c}}\}}$، أي:
\[
\Pow{\{ {\OLId{latin-ordinary}{a}}, {\OLId{latin-ordinary}{b}}, {\OLId{latin-ordinary}{c}} \}} = \{\emptyset, \{{\OLId{latin-ordinary}{a}} \}, \{{\OLId{latin-ordinary}{b}}\}, \{{\OLId{latin-ordinary}{c}}\}, \{{\OLId{latin-ordinary}{a}}, {\OLId{latin-ordinary}{b}}\},
\{{\OLId{latin-ordinary}{b}}, {\OLId{latin-ordinary}{c}}\}, \{{\OLId{latin-ordinary}{a}}, {\OLId{latin-ordinary}{c}}\}, \{{\OLId{latin-ordinary}{a}}, {\OLId{latin-ordinary}{b}}, {\OLId{latin-ordinary}{c}}\}\}
\]
\end{ex}

\begin{prob}
اذكر جزئيات $\{{\OLId{latin-ordinary}{a}}, {\OLId{latin-ordinary}{b}}, {\OLId{latin-ordinary}{c}}, {\OLId{latin-ordinary}{d}}\}$ جميعًا.
\end{prob}

\begin{prob}
بيّن أن ${\OLId{latin-looped}{A}}$ إذا كان عدد !!{element}sها ${\OLId{latin-ordinary}{n}}$، كان عدد !!{element}s
المجموعة $\Pow{{\OLId{latin-looped}{A}}}$ هو ${\OLMathNumeral{2}}^{\OLId{latin-ordinary}{n}}$.
\end{prob}

\end{document}
